\section{1 INTRODUÇÃO}\label{introduuxe7uxe3o}

\subsection{1.1 Contextualização}\label{contextualizauxe7uxe3o}

A educação inclusiva consolidou-se, nas últimas décadas, como um dos
principais paradigmas das políticas educacionais brasileiras. O
princípio da equidade e o direito universal à educação de qualidade
estão inscritos na Constituição Federal de 1988 e foram desdobrados pela
Lei de Diretrizes e Bases da Educação Nacional (LDB), de 1996, pela
Política Nacional de Educação Especial na Perspectiva da Educação
Inclusiva, de 2008, pela Lei 12.764/2012, que instituiu a Política
Nacional de Proteção dos Direitos da Pessoa com Transtorno do Espectro
Autista, e pela Lei Brasileira de Inclusão (LBI), de 2015. Esse
arcabouço estabelece que a escola comum é o espaço ordinário de
escolarização e que o atendimento educacional especializado deve
complementar, e não substituir, a experiência escolar regular (BRASIL,
1988; BRASIL, 1996; BRASIL, 2008; BRASIL, 2012; BRASIL, 2015).

A inclusão, contudo, não se resolve no plano normativo. A escola é uma
instituição atravessada por condições materiais, por concepções sobre
aprendizagem e por relações de poder que podem aproximar ou afastar a
presença legal da participação efetiva. No caso do Transtorno do
Espectro Autista (TEA), essa distância tende a se acentuar. O TEA é uma
condição do neurodesenvolvimento que se manifesta em um espectro amplo
de configurações: dificuldades na comunicação e na interação social,
padrões restritos e repetitivos de comportamento, interesses intensos e
respostas sensoriais atípicas podem se combinar de maneiras muito
distintas. Não há um aluno típico com TEA; há trajetórias singulares que
exigem planejamento pedagógico igualmente singular (CABRAL; MARIN, 2017;
NUNES; AZEVEDO; SCHMIDT, 2013).

Os anos iniciais do ensino fundamental constituem um período
especialmente sensível. É nessa fase que se consolidam bases da
alfabetização, do raciocínio matemático, das habilidades sociais e da
relação com a escola. Quando o ambiente não oferece previsibilidade,
apoio à comunicação e mediação pedagógica, o estudante pode acumular
experiências de fracasso que não decorrem de sua condição, mas da
inadequação das condições de ensino. A inclusão, portanto, não é apenas
uma questão de acesso: é uma questão de permanência, participação e
aprendizagem (CABRAL; MARIN, 2017; SCHMIDT et al., 2016).

A literatura educacional tem produzido diagnósticos consistentes sobre
os obstáculos enfrentados por esses estudantes. Muitos professores
relatam não se sentir preparados para lidar com as necessidades
específicas do TEA; faltam recursos pedagógicos, apoio especializado e
tempo institucional para planejamento colaborativo. Ao mesmo tempo,
persistem práticas que reduzem a inclusão à presença física do aluno na
sala regular, sem alteração do currículo, das formas de avaliação ou da
organização do trabalho pedagógico. O reconhecimento dessas barreiras,
porém, nem sempre vem acompanhado de uma análise de como políticas,
concepções e práticas se articulam --- ou se desarticulam --- no
cotidiano escolar (CABRAL; MARIN, 2017; COUTO et al., 2019; SCHMIDT et
al., 2016).

\subsection{1.2 Problema e pergunta de
pesquisa}\label{problema-e-pergunta-de-pesquisa}

O problema que orienta este trabalho pode ser formulado nos seguintes
termos: embora o direito à educação inclusiva esteja amplamente
reconhecido na legislação brasileira, as condições efetivas para a
participação e a aprendizagem de estudantes com TEA nos anos iniciais
ainda são desiguais e dependem de fatores que transcendem a matrícula. A
distância entre a norma e a prática não é um detalhe conjuntural; ela
expressa contradições entre o projeto inclusivo da escola e a
organização do trabalho pedagógico, da formação docente e da gestão
escolar.

Diante disso, a pergunta que orienta a investigação é: \textbf{como
políticas educacionais, práticas pedagógicas e formação docente se
articulam, na literatura nacional, para viabilizar a participação e a
aprendizagem de estudantes com Transtorno do Espectro Autista nos anos
iniciais do ensino fundamental?}

A pergunta não busca medir a eficácia de uma intervenção específica, nem
estimar a prevalência do TEA. Busca compreender como três dimensões ---
a política, a prática e a formação --- se relacionam na produção do que
se convencionou chamar de inclusão escolar. Trata-se, portanto, de uma
pergunta de natureza interpretativa, adequada a uma revisão narrativa da
literatura.

\subsection{1.3 Objetivos}\label{objetivos}

O objetivo geral deste trabalho é analisar como políticas educacionais,
práticas pedagógicas e formação docente se articulam, na literatura
nacional, para viabilizar a participação e a aprendizagem de estudantes
com TEA nos anos iniciais.

Como objetivos específicos, busca-se:

\begin{enumerate}
\def\labelenumi{\alph{enumi})}
\item
  situar o debate sobre as políticas de inclusão e seus limites de
  implementação;
\item
  discutir as implicações educacionais da variabilidade do TEA;
\item
  sistematizar práticas pedagógicas inclusivas voltadas aos anos
  iniciais;
\item
  identificar tensões transversais que atravessam a inclusão de
  estudantes com TEA;
\item
  propor uma agenda de pesquisa e de formação docente que contribua para
  o enfrentamento dessas tensões.
\end{enumerate}

\subsection{1.4 Justificativa}\label{justificativa}

A relevância do estudo se sustenta em três argumentos. O primeiro é de
natureza ética: a inclusão escolar é um direito, não uma concessão. A
presença de estudantes com TEA na escola comum não deve ser tratada como
favor ou como problema a ser administrado, mas como expressão do direito
à educação em igualdade de condições (BRASIL, 1988; BRASIL, 2015).

O segundo é de natureza prática: professores, gestores e famílias
precisam de referenciais que os ajudem a tomar decisões em contextos de
incerteza. A literatura sobre inclusão é extensa, mas frequentemente
fragmentada em estudos de caso, revisões parciais e relatos de
experiência. Uma síntese que organize os achados em eixos e tensões pode
facilitar o diálogo entre pesquisa e prática.

O terceiro é de natureza acadêmica: há lacunas na forma como a
literatura nacional articula políticas, concepções e práticas. Muitos
estudos tratam de um desses elementos isoladamente; poucos os colocam em
relação. Ao propor uma síntese integrativa, este trabalho pretende
contribuir para o campo da educação especial e para a formação inicial e
continuada de professores.

\subsection{1.5 Organização do
trabalho}\label{organizauxe7uxe3o-do-trabalho}

O trabalho está organizado em seis seções, além desta introdução. A
seção 2 apresenta a revisão da literatura, dividida em três eixos:
políticas educacionais de inclusão, concepções e implicações
educacionais do TEA e práticas pedagógicas inclusivas nos anos iniciais.
A seção 3 descreve a metodologia: tipo de pesquisa, corpus, critérios de
seleção, procedimentos de análise, verificação bibliográfica, limitações
e aspectos éticos. A seção 4 apresenta os resultados, com a
caracterização do corpus e a síntese dos achados por eixo. A seção 5
discute os resultados à luz de quatro tensões transversais e de suas
implicações para a formação docente e para a gestão escolar. A seção 6
apresenta as considerações finais. Em seguida, constam as referências.

\section{2 REVISÃO DA LITERATURA}\label{revisuxe3o-da-literatura}

\subsection{2.1 Políticas educacionais de inclusão no
Brasil}\label{poluxedticas-educacionais-de-inclusuxe3o-no-brasil}

O marco legal da educação inclusiva no Brasil é extenso e relativamente
consistente em seus princípios. A Constituição Federal de 1988
estabelece a educação como direito de todos e dever do Estado, com
atendimento educacional especializado aos portadores de deficiência,
preferencialmente na rede regular de ensino (BRASIL, 1988). A LDB de
1996 reafirma esse princípio e organiza a educação especial como
modalidade transversal (BRASIL, 1996). A Política Nacional de Educação
Especial na Perspectiva da Educação Inclusiva, de 2008, desloca o foco
do modelo segregacionista para a escola comum e institui o atendimento
educacional especializado como apoio complementar (BRASIL, 2008). A Lei
12.764/2012 reconhece a pessoa com TEA como pessoa com deficiência para
todos os efeitos legais e estabelece diretrizes de proteção (BRASIL,
2012). A LBI de 2015 amplia esse conjunto e reafirma o direito à
educação em igualdade de condições (BRASIL, 2015).

O Quadro 1 sintetiza os principais marcos e sua contribuição para a
inclusão de estudantes com TEA.

\textbf{Quadro 1 --- Marco legal da educação inclusiva no Brasil}

\begingroup\small

\begin{longtable}[]{@{}
  >{\raggedright\arraybackslash}p{(\linewidth - 4\tabcolsep) * \real{0.3333}}
  >{\raggedright\arraybackslash}p{(\linewidth - 4\tabcolsep) * \real{0.3333}}
  >{\raggedright\arraybackslash}p{(\linewidth - 4\tabcolsep) * \real{0.3333}}@{}}
\toprule\noalign{}
\begin{minipage}[b]{\linewidth}\raggedright
Norma
\end{minipage} & \begin{minipage}[b]{\linewidth}\raggedright
Ano
\end{minipage} & \begin{minipage}[b]{\linewidth}\raggedright
Contribuição principal
\end{minipage} \\
\midrule\noalign{}
\endhead
\bottomrule\noalign{}
\endlastfoot
Constituição Federal & 1988 & Educação como direito de todos;
atendimento educacional especializado preferencialmente na rede
regular \\
LDB 9.394 & 1996 & Educação especial como modalidade transversal;
organização da educação básica \\
PNEEPEI & 2008 & Escola comum como espaço ordinário; atendimento
educacional especializado complementar \\
Lei 12.764 & 2012 & Reconhecimento da pessoa com TEA como pessoa com
deficiência; diretrizes de proteção \\
Lei 13.146 (LBI) & 2015 & Direito à educação em igualdade de condições;
acessibilidade e não discriminação \\
\end{longtable}

\endgroup

Fonte: elaborado pela autora com base na legislação brasileira (BRASIL,
1988; BRASIL, 1996; BRASIL, 2008; BRASIL, 2012; BRASIL, 2015).

A existência desse arcabouço, no entanto, não garante sua efetivação. A
literatura analisada argumenta que a simples promulgação de leis é
insuficiente: é preciso que as políticas sejam acompanhadas de ações
práticas que assegurem condições de implementação (CABRAL; MARIN, 2017;
NUNES; AZEVEDO; SCHMIDT, 2013). Essa constatação desloca a análise do
plano jurídico para o plano organizacional. Não basta saber o que a lei
determina; é necessário perguntar quais recursos, quais profissionais,
quais tempos e quais espaços tornam a determinação realizável. A
inclusão, nesse sentido, é menos um enunciado de direito e mais um
conjunto de decisões cotidianas sobre currículo, avaliação, agrupamento
de estudantes e apoio pedagógico.

Nunes, Azevedo e Schmidt (2013), em revisão da literatura nacional sobre
inclusão educacional de pessoas com autismo, observam que a presença
desses estudantes em escolas regulares aumentou de forma expressiva após
a popularização do paradigma inclusivo. Esse crescimento, porém, não foi
acompanhado, na mesma proporção, de conhecimento sobre o TEA e de
estratégias pedagógicas específicas. O resultado é uma inclusão que os
autores descrevem como parcial: o estudante está na escola, mas nem
sempre participa das atividades com as mesmas possibilidades de
aprendizagem. A ampliação do acesso, portanto, não esgota o direito à
educação; ela apenas cria as condições para que o direito à participação
e à aprendizagem possa ser disputado.

Schmidt et al.~(2016), ao investigarem a percepção docente e as práticas
pedagógicas na inclusão de estudantes com autismo, encontram um quadro
ambivalente. De um lado, os professores reconhecem o direito à matrícula
e valorizam a convivência com a diversidade; de outro, relatam
insegurança, falta de formação e ausência de suporte para conduzir o
trabalho pedagógico. Essa ambivalência é importante porque não se reduz
a resistência individual. Ela expressa uma contradição institucional: a
escola é chamada a realizar uma tarefa para a qual não foi preparada, e
o professor é responsabilizado por um fracasso que tem origem em
condições coletivas. Responsabilizar apenas o docente, nesse cenário, é
uma forma de transferir o problema sem enfrentá-lo.

A implementação das políticas depende, ainda, de articulação
intersetorial. Cabral e Marin (2017) destacam a importância da parceria
entre família e escola e do trabalho interdisciplinar; a inclusão,
nesses termos, não é responsabilidade exclusiva da escola. Essa
perspectiva é especialmente relevante no caso do TEA, em que o apoio
especializado pode envolver fonoaudiologia, terapia ocupacional,
psicologia e acompanhamento pedagógico. A ausência de uma rede integrada
tende a produzir dois efeitos indesejados: a escola se isola diante de
demandas que não pode responder sozinha, e as famílias passam a circular
entre serviços que não se comunicam. A intersetorialidade, nesse
sentido, não é um adorno programático; é uma condição de possibilidade
da inclusão.

O monitoramento e a avaliação das políticas também merecem atenção. Fam
et al.~(2021) defendem que a avaliação das práticas inclusivas seja um
processo dinâmico e participativo, envolvendo educadores, estudantes e
famílias. Sem mecanismos de acompanhamento, a política tende a se tornar
um conjunto de metas declaradas, sem verificação de efeitos. A
avaliação, aqui, não deve ser confundida com ranqueamento de escolas ou
de professores; ela é um instrumento de aprendizagem institucional, que
permite identificar o que funciona, para quem e em quais condições. A
ausência dessa cultura avaliativa contribui para que os mesmos problemas
se reproduzam sem que haja acúmulo de conhecimento sobre soluções.

A formação docente atravessa todos esses elementos. Bezerra e Pantoni
(2022), ao discutirem a formação de professores para a inclusão de
estudantes com TEA no Ensino Médio Integrado, evidenciam que a lacuna
formativa não é exclusiva dos anos iniciais. Os autores argumentam que a
formação continuada precisa articular conhecimento sobre o TEA,
estratégias pedagógicas e reflexão sobre a prática. Apesar de o estudo
incidir sobre outro nível de ensino, sua lógica é pertinente ao problema
aqui analisado: a formação não se resolve com informações genéricas
sobre o espectro, mas com oportunidades de analisar situações concretas,
planejar adaptações e revisar decisões à luz dos resultados.

Em síntese, a literatura analisada permite distinguir três lacunas nas
políticas de inclusão. A primeira é \textbf{normativa}: embora o
arcabouço legal seja amplo, há tensões entre princípios e
regulamentações específicas. A segunda é \textbf{operacional}: faltam
recursos, profissionais, tempo e condições materiais para que a política
se realize. A terceira é \textbf{avaliativa}: inexistem, em muitos
contextos, mecanismos sistemáticos de acompanhamento dos efeitos das
práticas. Essas lacunas não são independentes; elas se reforçam
mutuamente. Uma política bem formulada, mas sem recursos, tende a
produzir práticas frágeis; práticas frágeis, sem avaliação, tendem a
perpetuar a ideia de que a inclusão é impossível.

\subsection{2.2 Transtorno do Espectro Autista: conceitos e implicações
educacionais}\label{transtorno-do-espectro-autista-conceitos-e-implicauxe7uxf5es-educacionais}

O Transtorno do Espectro Autista é uma condição do neurodesenvolvimento
cuja definição contemporânea se organiza em dois grandes domínios:
déficits persistentes na comunicação e na interação social em múltiplos
contextos; e padrões restritos e repetitivos de comportamento,
interesses ou atividades. A esses domínios associam-se, com frequência,
alterações no processamento sensorial e dificuldades de regulação. O
termo ``espectro'' não designa uma gradação linear de gravidade, mas a
ampla variabilidade de manifestações: duas pessoas com o mesmo
diagnóstico podem apresentar necessidades muito diferentes (CABRAL;
MARIN, 2017; NUNES; AZEVEDO; SCHMIDT, 2013).

Essa variabilidade tem consequências diretas para a educação. Um
planejamento pedagógico baseado em um suposto ``aluno autista típico''
tende a fracassar, porque não corresponde a nenhum estudante real. O que
a literatura sugere é o inverso: a necessidade de avaliação cuidadosa
das competências e das barreiras de cada estudante, seguida de
adaptações específicas. Nunes, Azevedo e Schmidt (2013) enfatizam que o
conhecimento aprofundado sobre as características do autismo é condição
para que os educadores identifiquem as estratégias mais adequadas. Esse
conhecimento, porém, não deve ser usado para rotular ou para fixar
expectativas; deve orientar decisões pedagógicas provisórias, sujeitas a
revisão.

O diagnóstico precoce é frequentemente apontado como fator favorável ao
desenvolvimento. Couto et al.~(2019), ao analisarem experiências de
professores com o autismo, relacionam o diagnóstico precoce à
possibilidade de intervenções mais tempestivas e à inclusão escolar. A
relação, contudo, não é automática. Um diagnóstico pode abrir acesso a
apoios, mas também pode produzir estigma, reduzir expectativas e
justificar práticas segregadoras. A questão educacional não é apenas
``como identificar'', mas ``o que fazer com a identificação''. A
resposta a essa segunda pergunta depende de concepções sobre
deficiência, aprendizagem e papel da escola.

Cabral e Marin (2017), em revisão sistemática sobre a inclusão escolar
de crianças com TEA, mostram que a escola precisa se reorganizar para
responder às necessidades desses estudantes, o que envolve mudanças na
interação dos profissionais, flexibilização curricular e apoio à
aprendizagem. As autoras analisaram 25 estudos nacionais e 92
internacionais e identificaram que a maioria das pesquisas é empírica e
qualitativa, com foco nas dificuldades de comunicação, no
desconhecimento das características do TEA e na carência de estratégias
pedagógicas. Esse panorama é importante porque desloca a explicação do
fracasso escolar do indivíduo para o ambiente: não se trata de ``ensinar
o aluno a se adaptar'', mas de transformar as condições de ensino.

Schmidt et al.~(2016), ao investigarem a percepção docente e as práticas
pedagógicas na inclusão de estudantes com autismo, mostram que os
professores tendem a reconhecer a importância da inclusão, mas relatam
dificuldades em traduzi-la em estratégias cotidianas. Essa distância
entre adesão ao princípio e realização da prática é um dos achados mais
recorrentes na literatura. Ela sugere que a formação não pode se limitar
a sensibilização ou a informações sobre o espectro. É necessário
trabalhar com situações concretas: como organizar a rotina, como adaptar
uma atividade, como mediar a interação com os colegas, como avaliar o
progresso sem comparar o estudante a um padrão inexistente.

O brincar ocupa lugar central nesse processo. Fam et al.~(2021)
argumentam que o brincar no espaço escolar pode funcionar como
estratégia de inclusão, pois cria oportunidades de interação, de
comunicação e de participação em atividades compartilhadas. O brincar,
porém, não é espontaneamente inclusivo. Ele exige mediação, estrutura e
adaptação. Sem apoio, a criança com TEA pode permanecer à margem das
brincadeiras, mesmo em um pátio cheio de colegas. A literatura, nesse
ponto, converge para uma ideia mais ampla: a inclusão se realiza em
situações sociais concretas, não em declarações de princípio.

A formação docente, novamente, aparece como elemento transversal. Cabral
e Marin (2017) e Schmidt et al.~(2016) destacam a necessidade de
programas de capacitação que ofereçam suporte prático aos professores. O
argumento não é que o professor deva dominar técnicas clínicas, mas que
precise compreender como o TEA se manifesta na sala de aula e como
ajustar o ensino em função disso. A formação, nesse sentido, é uma
prática de tradução: traduzir conhecimento sobre o desenvolvimento para
decisões pedagógicas situadas.

Em síntese, o TEA não é um conjunto de déficits a serem corrigidos, mas
uma forma de funcionamento que demanda condições educacionais
específicas. A literatura analisada sugere que a escola precisa combinar
três movimentos: reconhecer a variabilidade, organizar o ambiente para
reduzir barreiras e mediar a participação social. Nenhum desses
movimentos se realiza sem formação, sem tempo e sem apoio institucional.
A inclusão, portanto, não depende apenas de conhecer o diagnóstico;
depende de construir condições para que o estudante possa aprender e
conviver.

\subsection{2.3 Práticas pedagógicas inclusivas nos anos
iniciais}\label{pruxe1ticas-pedaguxf3gicas-inclusivas-nos-anos-iniciais}

As práticas pedagógicas inclusivas não constituem um receituário. Elas
são respostas situadas a barreiras identificadas em contextos
específicos. A literatura analisada, no entanto, permite agrupar um
conjunto de estratégias que aparecem de forma recorrente: adaptação
curricular, planejamento diversificado, ensino estruturado, comunicação
alternativa, tecnologias assistivas, mediação do brincar, trabalho
colaborativo e participação da família. Essas estratégias não são
equivalentes nem intercambiáveis; sua pertinência depende do estudante,
do momento e das condições disponíveis (CABRAL; MARIN, 2017; NUNES,
2018).

Cabral e Marin (2017) e Schmidt et al.~(2016) problematizam a relação
entre os apoios disponíveis e as possibilidades de aprendizagem e
desenvolvimento. A literatura reforça a necessidade de examinar as
condições de ensino, e não apenas o acesso, quando se avalia a inclusão.

A adaptação curricular é frequentemente apresentada como eixo central.
Cabral e Marin (2017) e Bezerra e Pantoni (2022) destacam que a
adaptação deve considerar as particularidades de cada aluno e promover
sua participação ativa nas atividades. Adaptar, porém, não é simplificar
o currículo até o ponto de retirar do estudante a oportunidade de
aprender. A adaptação exige decisões sobre o que é essencial, quais são
os objetivos prioritários e quais caminhos permitem alcançá-los. Em vez
de reduzir expectativas, trata-se de diversificar percursos. Bezerra e
Pantoni (2022) reforçam que a formação continuada precisa preparar o
professor para tomar essas decisões, articulando conhecimento pedagógico
e conhecimento sobre o TEA.

O planejamento e a diversificação metodológica são condições para que a
adaptação se realize. Schmidt et al.~(2016) sugerem que o uso de
recursos visuais, auditivos e táteis pode facilitar a compreensão do
conteúdo por estudantes com autismo. A combinação de diferentes
linguagens amplia as possibilidades de acesso ao conhecimento e reduz a
dependência de uma única forma de expressão. O planejamento, nesse
sentido, não é um documento burocrático; é o espaço em que o professor
antecipa barreiras e imagina alternativas. Quando o planejamento é
padronizado, a adaptação tende a ocorrer apenas depois do fracasso, como
remediação.

O ensino estruturado e a Análise do Comportamento Aplicada (ABA)
aparecem na literatura como abordagens com resultados relatados na
promoção da aprendizagem e de habilidades sociais. Couto et al.~(2019)
apontam que essas estratégias têm sido utilizadas com estudantes com
TEA. É importante evitar dois extremos: trata-se de abordagens que podem
oferecer recursos importantes, mas não são a única via e não devem ser
aplicadas de forma mecânica, sem consideração ao contexto e à
singularidade do estudante. A eficácia relatada na literatura não
dispensa julgamento pedagógico nem acompanhamento sistemático dos
efeitos.

A comunicação alternativa merece destaque. Togashi e Walter (2016)
analisam as contribuições do uso da comunicação alternativa no processo
de inclusão escolar de um aluno com TEA e mostram que recursos de
comunicação podem ampliar a participação e a expressão. A comunicação,
aqui, não é um pré-requisito para a inclusão; é um direito que a escola
deve garantir. Quando a escola exige que o estudante se comunique de uma
única forma para participar, ela transforma uma barreira em critério de
exclusão. A comunicação alternativa inverte essa lógica: adapta o
ambiente para que a expressão seja possível.

As tecnologias assistivas ampliam esse princípio. Maia, Bataglion e Mazo
(2020), ao analisarem relatos de professores de educação física sobre
estudantes com TEA na escola regular, destacam que recursos e adaptações
podem favorecer a participação nas atividades. A tecnologia assistiva
não se limita a dispositivos sofisticados; inclui desde apoios visuais e
pranchas de comunicação até adaptações simples de materiais e do espaço.
O critério para sua adoção não é a novidade tecnológica, mas a redução
de uma barreira concreta. Uma solução simples, bem ajustada, pode ser
mais eficaz do que um recurso avançado que não dialoga com a rotina da
sala.

O brincar e as atividades colaborativas ocupam lugar central nos anos
iniciais. Fam et al.~(2021) defendem que o brincar no espaço escolar
pode ser uma estratégia de inclusão, pois cria situações de interação e
de aprendizagem compartilhada. Cabral e Marin (2017) e Fam et al.~(2021)
também destacam a importância de atividades lúdicas e colaborativas para
o desenvolvimento social. O brincar, contudo, precisa ser planejado e
mediado. A simples justaposição de crianças não garante interação; é
necessário estruturar papéis, oferecer apoios e ensinar habilidades
sociais que muitos estudantes dominam implicitamente, mas que outros
precisam aprender de forma explícita.

A avaliação é outro elemento decisivo. Cabral e Marin (2017) e Fam et
al.~(2021) destacam que a avaliação seja contínua, dinâmica e
participativa, envolvendo educadores, estudantes e famílias. Avaliar,
nesse contexto, não é medir o quanto o estudante se aproxima de um
padrão; é acompanhar o progresso em relação às suas próprias
possibilidades e às metas definidas no plano educacional. Uma avaliação
padronizada, aplicada sem adaptações, tende a produzir resultados que
refletem mais as barreiras do instrumento do que as competências do
estudante. A avaliação inclusiva exige, portanto, diversidade de formas
de expressão e de registro.

O trabalho colaborativo entre professores e profissionais de apoio é
condição para que essas práticas se sustentem. Cabral e Marin (2017)
argumentam que a abordagem interdisciplinar é essencial, pois permite
que diferentes profissionais contribuam com seus conhecimentos. A
colaboração, porém, não se decreta. Ela exige tempo institucional,
canais de comunicação e clareza sobre responsabilidades. Quando cada
profissional atua de forma isolada, o estudante recebe intervenções
fragmentadas, que podem se contradizer. A colaboração, nesse sentido, é
uma prática pedagógica em si, não um acessório.

A relação família-escola completa esse conjunto. Cabral e Marin (2017) e
Fam et al.~(2021) destacam que o envolvimento das famílias fortalece a
inclusão e contribui para a compreensão das necessidades e
potencialidades dos estudantes. Essa relação, no entanto, não deve ser
concebida como transferência de responsabilidade. A família não é uma
extensão da escola nem uma fonte de informação a ser extraída; é
parceira na tomada de decisões. O diálogo exige escuta, respeito às
diferenças culturais e reconhecimento de que o conhecimento da família
sobre o estudante é insubstituível.

A acessibilidade física também integra o conjunto de condições. Maia,
Bataglion e Mazo (2020) destacam que a organização do espaço e as
adaptações devem facilitar a participação de todos os alunos. No caso do
TEA, a acessibilidade não se reduz a rampas ou portas largas; envolve
sinalização, previsibilidade, controle de estímulos sensoriais e espaços
de regulação. Um ambiente fisicamente acessível, mas sensorialmente
hostil, continua sendo excludente. A concepção de acessibilidade
precisa, portanto, ser ampliada para incluir dimensões comunicacionais,
sensoriais e cognitivas.

Em síntese, as práticas pedagógicas inclusivas nos anos iniciais
combinam adaptação, mediação e colaboração. Nenhuma delas funciona
isoladamente. A adaptação curricular sem formação tende a se tornar
improviso; a formação sem tempo e recursos tende a se tornar discurso;
os recursos sem mediação tendem a se tornar objetos sem uso. A
literatura analisada sugere que a qualidade da inclusão depende menos da
adoção de uma técnica específica e mais da coerência entre objetivos,
condições e decisões pedagógicas.

\section{3 METODOLOGIA}\label{metodologia}

\subsection{3.1 Tipo e abordagem}\label{tipo-e-abordagem}

Este trabalho adota a revisão narrativa como procedimento de
investigação. A revisão narrativa é adequada quando o objetivo é mapear
conceitos, identificar lacunas e construir uma interpretação
fundamentada sobre um campo, e não estimar a magnitude de um efeito ou
esgotar a evidência disponível. Diferentemente da revisão sistemática,
ela não pressupõe protocolo fechado de busca, triagem em dupla ou
avaliação padronizada de risco de viés. Essa escolha é coerente com a
natureza do problema: trata-se de compreender como três dimensões ---
políticas, concepções e práticas --- se articulam em um campo
atravessado por disputas conceituais e por condições institucionais
heterogêneas.

A abordagem é qualitativa e interpretativa. O material analisado não é
tratado como um conjunto de dados a serem medidos, mas como um corpus de
argumentos, experiências e perspectivas que precisam ser lidos em sua
singularidade e em suas relações. A análise busca identificar
recorrências, tensões e lacunas, sem pretender produzir generalizações
estatísticas. O critério de rigor, aqui, é a coerência da interpretação
com as fontes, a transparência do percurso e a explicitação das
limitações.

\subsection{3.2 Corpus da revisão}\label{corpus-da-revisuxe3o}

O corpus reúne oito artigos publicados em periódicos, entre 2013 e 2022,
em português, sobre inclusão escolar de estudantes com TEA na educação
básica, com ênfase nos anos iniciais. Todos possuem DOI ativo e
auditável. A exigência de DOI foi adotada como critério de integridade e
rastreabilidade da literatura científica. A seleção partiu de um
conjunto de referências previamente reunido pela autora em sua
trajetória de leitura sobre o tema e foi orientada por critérios de
aderência temática, disponibilidade do texto e contribuição para os
eixos de análise.

Além das oito obras científicas com DOI ativo, o corpus inclui cinco
documentos legais que estruturam o debate: a Constituição Federal de
1988, a LDB de 1996, a PNEEPEI de 2008, a Lei 12.764/2012 e a LBI de
2015. Esses documentos não são tratados como literatura empírica, mas
como fontes primárias do ordenamento jurídico que define o direito à
educação inclusiva.

\subsection{3.3 Critérios de
seleção}\label{crituxe9rios-de-seleuxe7uxe3o}

Foram adotados os seguintes critérios de inclusão:

\begin{enumerate}
\def\labelenumi{\alph{enumi})}
\item
  abordagem simultânea do TEA e da inclusão escolar;
\item
  pertinência à educação básica, especialmente aos anos iniciais;
\item
  disponibilidade do texto em português;
\item
  contribuição para ao menos um dos eixos de análise: políticas,
  concepções sobre o TEA ou práticas pedagógicas.
\end{enumerate}

Foram excluídos:

\begin{enumerate}
\def\labelenumi{\alph{enumi})}
\item
  estudos estritamente clínicos, sem implicações educacionais;
\item
  obras cuja discussão sobre inclusão não incidisse sobre o TEA;
\item
  textos sem identificação de autoria ou sem possibilidade de
  verificação bibliográfica;
\item
  duplicatas ou versões repetidas de um mesmo trabalho.
\end{enumerate}

A seleção não foi exaustiva. O corpus é um recorte interpretativo, e não
uma amostra representativa de toda a produção nacional. Essa limitação é
assumida e discutida na seção 3.6.

\subsection{3.4 Procedimentos de
análise}\label{procedimentos-de-anuxe1lise}

A análise seguiu três etapas. Na primeira, realizou-se leitura integral
de cada obra, com registro de dados bibliográficos, objetivos,
procedimentos e principais argumentos. Na segunda, os conteúdos foram
categorizados em cinco núcleos: marco legal e político; concepções sobre
o TEA; estratégias pedagógicas; formação docente; e relação
família-escola. Na terceira, os núcleos foram reagrupados em três eixos
de exposição --- políticas, TEA e práticas --- e em uma síntese
integrativa, na qual se buscou identificar tensões transversais aos
materiais.

A categorização não foi realizada com software de análise qualitativa.
Foi um processo manual, orientado pela leitura repetida e pela
comparação entre os textos. Essa escolha tem limites: não há medida de
concordância entre codificadores, nem auditoria externa do processo. Em
contrapartida, permite preservar a singularidade dos argumentos e evitar
a fragmentação excessiva do material em códigos.

\subsection{3.5 Verificação
bibliográfica}\label{verificauxe7uxe3o-bibliogruxe1fica}

Como procedimento de verificação, os DOIs das referências científicas
foram resolvidos e conferidos na API Crossref em 5 de novembro de 2026.
Oito obras tinham DOI ativo e foram aceitas. Seis obras foram excluídas:
Costa (2016), Nascimento e Cruz (2015), Nunes (2018), Ribeiro (2019),
Santos e Lopes (2021) e Silva e Arantes (2020). A obra de Nascimento e
Cruz (2015) apresentava DOI impresso, mas não registrado (HTTP 404); as
demais não tinham DOI localizado. A obra de Schmidt et al.~(2016) tem
DOI registrado no Crossref; a página do editor bloqueia acesso
automatizado, mas o registro é auditável. A legislação é tratada como
fonte primária, com URL oficial, por não possuir DOI.

A verificação de DOI não equivale à leitura integral nem à avaliação da
qualidade das obras. Trata-se de um controle de integridade e
rastreabilidade bibliográfica, não de um julgamento de mérito. O
registro da auditoria permanece com a autora, disponível para
verificação.

\subsection{3.6 Limitações}\label{limitauxe7uxf5es}

O estudo apresenta limitações que devem ser consideradas na leitura dos
resultados. Por ser revisão narrativa, não há garantia de exaustividade;
a seleção do corpus pode privilegiar determinadas perspectivas em
detrimento de outras. Não foi realizada avaliação padronizada de risco
de viés, nem análise meta-analítica. O corpus é restrito à literatura
nacional e a um recorte temporal específico. Por isso, as conclusões são
interpretativas e não devem ser generalizadas como descrição de todas as
escolas brasileiras. Elas indicam tendências e tensões que a literatura
permite sustentar, não regularidades empíricas.

Há também limitações de natureza bibliográfica. A exigência de DOI ativo
reduziu o corpus de 14 para oito obras. Seis obras foram excluídas, e
suas contribuições foram absorvidas, quando possível, pelas obras
DOI-bearing remanescentes. A legislação é citada como fonte primária,
com URL oficial, mas não possui DOI. Essa restrição diminui a
diversidade do corpus e deve ser considerada na leitura dos resultados.
As decisões estão registradas no relatório de auditoria que acompanha
este trabalho.

\subsection{3.7 Aspectos éticos}\label{aspectos-uxe9ticos}

A pesquisa não envolveu seres humanos, coleta de dados primários ou
intervenção em contextos escolares. Todos os materiais utilizados são de
acesso público e foram citados com identificação de autoria. A análise
buscou respeitar a integridade das fontes, evitando distorções, citações
descontextualizadas ou atribuições indevidas. As obras não localizadas
ou com autoria divergente foram sinalizadas, em vez de silenciadas.

\section{4 RESULTADOS}\label{resultados}

\subsection{4.1 Caracterização do
corpus}\label{caracterizauxe7uxe3o-do-corpus}

O corpus científico é composto por oito artigos publicados em
periódicos, entre 2013 e 2022. Todos foram selecionados por atenderem ao
critério de DOI ativo e auditável. O Quadro 2 apresenta a caracterização
geral.

\textbf{Quadro 2 --- Caracterização do corpus científico}

\begingroup\small

\begin{longtable}[]{@{}
  >{\raggedright\arraybackslash}p{(\linewidth - 6\tabcolsep) * \real{0.2500}}
  >{\raggedright\arraybackslash}p{(\linewidth - 6\tabcolsep) * \real{0.2500}}
  >{\raggedright\arraybackslash}p{(\linewidth - 6\tabcolsep) * \real{0.2500}}
  >{\raggedright\arraybackslash}p{(\linewidth - 6\tabcolsep) * \real{0.2500}}@{}}
\toprule\noalign{}
\begin{minipage}[b]{\linewidth}\raggedright
Autor/ano
\end{minipage} & \begin{minipage}[b]{\linewidth}\raggedright
Tipo
\end{minipage} & \begin{minipage}[b]{\linewidth}\raggedright
Veículo
\end{minipage} & \begin{minipage}[b]{\linewidth}\raggedright
Contribuição principal
\end{minipage} \\
\midrule\noalign{}
\endhead
\bottomrule\noalign{}
\endlastfoot
Nunes, Azevedo e Schmidt (2013) & Artigo & Revista Educação Especial &
Revisão da literatura sobre inclusão de pessoas com autismo no Brasil \\
Schmidt et al.~(2016) & Artigo & Psicologia: Teoria e Prática &
Percepção docente e práticas pedagógicas na inclusão \\
Togashi e Walter (2016) & Artigo & Revista Brasileira de Educação
Especial & Comunicação alternativa na inclusão escolar \\
Cabral e Marin (2017) & Artigo & Educação em Revista & Revisão
sistemática da literatura nacional e internacional \\
Couto et al.~(2019) & Artigo & Revista Eletrônica de Enfermagem &
Experiências de professores com o autismo \\
Maia, Bataglion e Mazo (2020) & Artigo & Revista da ABAM & Relatos de
professores de educação física \\
Fam et al.~(2021) & Artigo & Research, Society and Development & O
brincar como estratégia de inclusão \\
Bezerra e Pantoni (2022) & Artigo & Educitec & Formação docente para
inclusão no Ensino Médio Integrado \\
\end{longtable}

\endgroup

Fonte: elaborado pela autora com base no corpus da pesquisa
(2013--2022).

A distribuição temporal é irregular: há uma obra de 2013, duas de 2016,
uma de 2017, uma de 2019, uma de 2020, uma de 2021 e uma de 2022. Essa
distribuição sugere que a inclusão de estudantes com TEA se consolidou
como objeto de pesquisa educacional sobretudo após a LBI de 2015, embora
o debate já estivesse presente na literatura anterior. A totalidade do
corpus é composta por artigos em periódicos, o que indica que o campo se
organiza em torno de revistas especializadas.

\subsection{4.2 Achados sobre políticas de
inclusão}\label{achados-sobre-poluxedticas-de-inclusuxe3o}

Os estudos convergem para a ideia de que o avanço normativo não foi
acompanhado, na mesma medida, de condições operacionais. Nunes, Azevedo
e Schmidt (2013) observam que o aumento da matrícula não foi acompanhado
de conhecimento sobre o TEA e de estratégias pedagógicas; Schmidt et
al.~(2016) mostram que os professores reconhecem o direito, mas relatam
insegurança e falta de suporte. A literatura analisada argumenta que a
promulgação de leis é insuficiente sem ações práticas (CABRAL; MARIN,
2017). Essa tríade --- norma, acesso e prática --- organiza a maior
parte dos achados sobre políticas.

Um segundo achado diz respeito à intersetorialidade. Cabral e Marin
(2017) destacam a importância da parceria entre família e escola e do
trabalho interdisciplinar. A ausência dessa rede produz isolamento da
escola e fragmentação do atendimento. Fam et al.~(2021) acrescentam a
necessidade de avaliação participativa das práticas, envolvendo
educadores, estudantes e famílias. A política, nesses termos, não se
esgota na formulação; depende de arranjos de implementação e de
monitoramento.

Um terceiro achado concerne à formação docente. Bezerra e Pantoni (2022)
mostram que a lacuna formativa não é exclusiva dos anos iniciais; ela
atravessa níveis de ensino e exige formação continuada articulada à
prática. Schmidt et al.~(2016) e Bezerra e Pantoni (2022) reforçam que o
professor precisa de tempo, apoio e oportunidades de reflexão para
transformar o princípio inclusivo em decisão pedagógica. A formação,
portanto, não é um apêndice da política; é uma de suas condições de
efetividade.

\subsection{4.3 Achados sobre o TEA e implicações
educacionais}\label{achados-sobre-o-tea-e-implicauxe7uxf5es-educacionais}

Os estudos sobre o TEA convergem para a centralidade da variabilidade.
Cabral e Marin (2017) mostram que a literatura nacional e internacional
enfatiza as dificuldades de comunicação, o desconhecimento das
características do TEA e a carência de estratégias pedagógicas. Nunes,
Azevedo e Schmidt (2013) destacam a necessidade de conhecimento
aprofundado para a identificação de estratégias adequadas. Schmidt et
al.~(2016) e Couto et al.~(2019) mostram que os professores reconhecem a
importância da inclusão, mas têm dificuldade em traduzi-la em práticas
cotidianas.

Outro achado é a relação entre diagnóstico e inclusão. Couto et
al.~(2019) relacionam o diagnóstico precoce à possibilidade de
intervenções tempestivas. Ao mesmo tempo, a literatura sugere cautela
para que a identificação não reduza o estudante à sua condição. O
diagnóstico pode abrir acesso a apoios, mas também pode produzir estigma
e reduzir expectativas. A questão educacional, portanto, não é apenas
identificar; é construir condições para que a identificação se converta
em apoio, e não em rótulo.

Um terceiro achado concerne ao papel do ambiente. Cabral e Marin (2017)
e Fam et al.~(2021) mostram que a inclusão depende da reorganização do
espaço, da previsibilidade das rotinas, das formas de comunicação e das
oportunidades de participação. Cabral e Marin (2017) problematizam a
relação entre apoios disponíveis e possibilidades de aprendizagem. A
inclusão, nesses termos, não é um atributo do estudante, mas uma
qualidade das condições de ensino.

\subsection{4.4 Achados sobre práticas
pedagógicas}\label{achados-sobre-pruxe1ticas-pedaguxf3gicas}

As práticas pedagógicas aparecem na literatura como um repertório
diversificado, e não como um método único. Cabral e Marin (2017) e
Bezerra e Pantoni (2022) destacam a adaptação curricular; Schmidt et
al.~(2016) enfatizam a diversificação de recursos visuais, auditivos e
táteis; Couto et al.~(2019) mencionam o ensino estruturado e a ABA;
Togashi e Walter (2016) evidenciam a comunicação alternativa; Maia,
Bataglion e Mazo (2020) tratam das tecnologias assistivas e das
adaptações na educação física; Fam et al.~(2021) destaca o brincar e as
atividades colaborativas; Cabral e Marin (2017) e Fam et al.~(2021)
tratam da avaliação contínua, do trabalho colaborativo e da relação
família-escola; Maia, Bataglion e Mazo (2020) abordam a acessibilidade e
as adaptações.

Um achado transversal é que nenhuma estratégia funciona isoladamente. A
adaptação curricular sem formação tende a se tornar improviso; a
formação sem tempo e recursos tende a se tornar discurso; os recursos
sem mediação tendem a se tornar objetos sem uso. A literatura sugere que
a qualidade da inclusão depende da coerência entre objetivos, condições
e decisões pedagógicas, e não da adoção de uma técnica específica.

Outro achado é a importância da mediação. O brincar, a comunicação
alternativa, as tecnologias assistivas e as atividades colaborativas não
são inclusivos por si mesmos. Eles se tornam inclusivos quando são
planejados, estruturados e acompanhados. A simples presença do recurso
ou da atividade não garante participação. A mediação, nesse sentido, é o
elemento que transforma um instrumento em oportunidade de aprendizagem.

\subsection{4.5 Síntese dos achados}\label{suxedntese-dos-achados}

A análise dos três eixos permite identificar seis padrões recorrentes:

\begin{enumerate}
\def\labelenumi{\arabic{enumi}.}
\tightlist
\item
  \textbf{Avanço normativo sem condições operacionais}: a lei garante o
  direito, mas a efetivação depende de recursos, formação e gestão.
\item
  \textbf{Formação docente como gargalo}: a formação inicial é
  insuficiente; a continuada, quando existe, é fragmentada e distante da
  prática.
\item
  \textbf{Escassez de recursos e apoio especializado}: faltam materiais,
  profissionais e tempo institucional para o planejamento colaborativo.
\item
  \textbf{Rigidez curricular e avaliativa}: o currículo e a avaliação
  padronizados dificultam a adaptação e a participação.
\item
  \textbf{Barreiras atitudinais}: preconceito, baixa expectativa e
  resistência à diversidade persistem, mesmo entre profissionais que
  aderem ao princípio inclusivo.
\item
  \textbf{Estratégias promissoras, mas condicionadas}: adaptação
  curricular, comunicação alternativa, tecnologias assistivas, brincar
  mediado, trabalho colaborativo e participação da família aparecem como
  promissoras, desde que inseridas em condições institucionais
  adequadas.
\end{enumerate}

\textbf{Quadro 3 --- Síntese dos achados por eixo}

\begingroup\small

\begin{longtable}[]{@{}
  >{\raggedright\arraybackslash}p{(\linewidth - 4\tabcolsep) * \real{0.3333}}
  >{\raggedright\arraybackslash}p{(\linewidth - 4\tabcolsep) * \real{0.3333}}
  >{\raggedright\arraybackslash}p{(\linewidth - 4\tabcolsep) * \real{0.3333}}@{}}
\toprule\noalign{}
\begin{minipage}[b]{\linewidth}\raggedright
Eixo
\end{minipage} & \begin{minipage}[b]{\linewidth}\raggedright
Achados recorrentes
\end{minipage} & \begin{minipage}[b]{\linewidth}\raggedright
Implicação
\end{minipage} \\
\midrule\noalign{}
\endhead
\bottomrule\noalign{}
\endlastfoot
Políticas & Norma forte, implementação frágil; intersetorialidade
incipiente; avaliação ausente & Política de implementação e
monitoramento \\
TEA & Variabilidade; diagnóstico ambivalente; ambiente como barreira ou
apoio & Conhecimento do TEA e reorganização do ensino \\
Práticas & Repertório diversificado; mediação necessária; colaboração
condicionante & Formação prática e tempo institucional \\
\end{longtable}

\endgroup

Fonte: elaborado pela autora com base no corpus da pesquisa.

\section{5 DISCUSSÃO}\label{discussuxe3o}

\subsection{5.1 Quatro tensões da
inclusão}\label{quatro-tensuxf5es-da-inclusuxe3o}

A leitura conjunta dos resultados permite identificar quatro tensões que
atravessam a literatura e a prática escolar. Essas tensões não são
problemas a serem eliminados, mas contradições a serem administradas.
Nomeá-las pode ajudar gestores, professores e pesquisadores a evitar
soluções simplistas.

\textbf{Tensão 1 --- Norma e prática.} O arcabouço legal brasileiro é
amplo e relativamente avançado em seus princípios, mas a efetivação
depende de condições que não são garantidas pela lei. Nunes, Azevedo e
Schmidt (2013) e Cabral e Marin (2017) mostram que a ampliação do acesso
não foi acompanhada, na mesma medida, de formação, recursos e
estratégias. A consequência é uma inclusão formal, que pode coexistir
com exclusão pedagógica. A tensão entre norma e prática não se resolve
com mais leis; resolve-se com políticas de implementação, financiamento
e acompanhamento.

\textbf{Tensão 2 --- Identificação e estigma.} O diagnóstico pode abrir
acesso a apoios e orientar intervenções, mas também pode reduzir o
estudante à sua condição. Couto et al.~(2019) relacionam o diagnóstico
precoce à possibilidade de intervenções tempestivas; a literatura
analisada sugere cautela para que a identificação não reduza o estudante
à sua condição. A questão é como usar a identificação sem fixar
expectativas. Uma resposta possível é deslocar o foco do diagnóstico
para as barreiras: em vez de perguntar ``o que falta ao estudante'',
perguntar ``o que o ambiente precisa oferecer para que ele participe''.
Esse deslocamento não nega a condição; ele redefine o objeto da
intervenção.

\textbf{Tensão 3 --- Adaptação e segregação.} Adaptar o currículo é
necessário, mas a adaptação pode se converter em redução de expectativas
ou em atividades paralelas que isolam o estudante. Cabral e Marin (2017)
e Bezerra e Pantoni (2022) defendem adaptações que promovam participação
ativa, não exclusão disfarçada. A diferença é sutil e decisiva: uma
adaptação inclusiva amplia o acesso ao conhecimento comum; uma adaptação
segregadora cria um currículo à parte, que não dialoga com o currículo
da turma. O critério para distinguir uma da outra é a participação: o
estudante está aprendendo o que os colegas aprendem, com apoios
adequados, ou está realizando outra coisa, em outro espaço, sem relação
com o projeto pedagógico coletivo?

\textbf{Tensão 4 --- Formação inicial e formação continuada.} A formação
inicial frequentemente não aborda, com profundidade, as necessidades
educacionais específicas do TEA. A formação continuada pode suprir parte
dessa lacuna, mas não resolve o problema estrutural. Bezerra e Pantoni
(2022) e Cabral e Marin (2017) argumentam que a formação precisa ser
prática, articulada ao cotidiano e acompanhada de suporte institucional.
A tensão, aqui, é entre responsabilizar o professor individualmente e
construir condições coletivas de formação. Sem tempo, sem apoio e sem
continuidade, a formação continuada se reduz a eventos isolados, que não
alteram práticas.

\textbf{Quadro 4 --- Quatro tensões da inclusão}

\begingroup\small

\begin{longtable}[]{@{}
  >{\raggedright\arraybackslash}p{(\linewidth - 6\tabcolsep) * \real{0.2500}}
  >{\raggedright\arraybackslash}p{(\linewidth - 6\tabcolsep) * \real{0.2500}}
  >{\raggedright\arraybackslash}p{(\linewidth - 6\tabcolsep) * \real{0.2500}}
  >{\raggedright\arraybackslash}p{(\linewidth - 6\tabcolsep) * \real{0.2500}}@{}}
\toprule\noalign{}
\begin{minipage}[b]{\linewidth}\raggedright
Tensão
\end{minipage} & \begin{minipage}[b]{\linewidth}\raggedright
Manifestação
\end{minipage} & \begin{minipage}[b]{\linewidth}\raggedright
Risco
\end{minipage} & \begin{minipage}[b]{\linewidth}\raggedright
Princípio orientador
\end{minipage} \\
\midrule\noalign{}
\endhead
\bottomrule\noalign{}
\endlastfoot
Norma e prática & Lei ampla, implementação frágil & Inclusão formal sem
aprendizagem & Política de implementação e monitoramento \\
Identificação e estigma & Diagnóstico como apoio ou rótulo & Redução de
expectativas & Foco nas barreiras do ambiente \\
Adaptação e segregação & Currículo adaptado ou paralelo & Exclusão
disfarçada & Participação no currículo comum \\
Formação inicial e continuada & Lacuna formativa e eventos isolados &
Culpabilização do professor & Formação prática e institucional \\
\end{longtable}

\endgroup

Fonte: elaborado pela autora com base na análise do corpus.

\subsection{5.2 Implicações para a formação
docente}\label{implicauxe7uxf5es-para-a-formauxe7uxe3o-docente}

As tensões identificadas têm implicações diretas para a formação inicial
e continuada. A primeira é que a formação não pode se limitar a
conteúdos sobre o TEA. É necessário trabalhar com situações concretas de
sala de aula: como organizar a rotina, como adaptar uma atividade, como
mediar a interação entre pares, como avaliar o progresso sem comparar o
estudante a um padrão inexistente. A formação, nesse sentido, precisa
ser prática e situada.

A segunda implicação é que a formação continuada deve ser acompanhada de
suporte institucional. Sem tempo para planejamento, sem materiais e sem
apoio especializado, o professor tende a voltar às práticas anteriores.
A formação, isolada, não transforma condições. Ela precisa estar
articulada a uma política de gestão que reconheça o tempo de
planejamento como parte do trabalho docente.

A terceira implicação é que a formação deve incluir a dimensão
colaborativa. O professor não pode ser o único responsável pela
inclusão. A colaboração com outros professores, com profissionais de
apoio, com a família e com a rede de saúde é condição para que o
trabalho se sustente. A formação, portanto, precisa desenvolver
competências de trabalho em equipe, de comunicação e de mediação de
conflitos.

\subsection{5.3 Implicações para a gestão
escolar}\label{implicauxe7uxf5es-para-a-gestuxe3o-escolar}

A gestão escolar tem um papel decisivo na transformação das condições.
Isso inclui a organização do espaço, a previsibilidade das rotinas, a
alocação de recursos, a distribuição do tempo e a criação de canais de
comunicação com as famílias. A gestão também é responsável por construir
uma cultura escolar que valorize a diversidade e enfrente as barreiras
atitudinais.

Uma gestão inclusiva não se limita a cumprir a matrícula. Ela acompanha
a participação, avalia os efeitos das práticas e ajusta o planejamento.
Ela reconhece que a inclusão é um processo coletivo, e não uma tarefa
individual do professor. Ela cria condições para que a formação
continuada se realize e para que o trabalho colaborativo seja possível.

\subsection{5.4 Limitações do estudo}\label{limitauxe7uxf5es-do-estudo}

As limitações já foram anunciadas na metodologia e devem ser retomadas
aqui. A primeira é a natureza narrativa da revisão: não há garantia de
exaustividade, nem avaliação padronizada de risco de viés. A segunda é a
restrição do corpus à literatura nacional e a um recorte temporal
específico. A terceira é a exclusão de seis obras que não atendiam à
regra de DOI ativo, o que reduziu a diversidade do corpus e exige
cautela na generalização.

Essas limitações não anulam os achados, mas delimitam seu alcance. O
estudo não pretende descrever todas as escolas brasileiras, nem estimar
a eficácia de intervenções. Pretende oferecer uma síntese interpretativa
que ajude a compreender as tensões e a orientar novas pesquisas.

\subsection{5.5 Agenda de pesquisa}\label{agenda-de-pesquisa}

A agenda de pesquisa que se abre a partir deste trabalho inclui:

\begin{enumerate}
\def\labelenumi{\alph{enumi})}
\item
  estudos empíricos que observem práticas de sala de aula em contextos
  diversos, com atenção aos efeitos das adaptações sobre a aprendizagem
  e a participação;
\item
  pesquisas longitudinais que acompanhem trajetórias de estudantes com
  TEA nos anos iniciais e identifiquem fatores que favorecem ou
  dificultam a permanência;
\item
  estudos sobre formação docente que avaliem programas articulados à
  prática, em vez de cursos pontuais;
\item
  pesquisas que incorporem a voz dos estudantes e de suas famílias como
  fonte de conhecimento sobre as barreiras e as possibilidades da
  inclusão;
\item
  estudos comparativos sobre políticas de implementação em diferentes
  redes de ensino.
\end{enumerate}

\section{6 CONCLUSÃO}\label{conclusuxe3o}

A pergunta que orientou este trabalho foi como políticas educacionais,
práticas pedagógicas e formação docente se articulam, na literatura
nacional, para viabilizar a participação e a aprendizagem de estudantes
com TEA nos anos iniciais. A análise permite responder que essa
articulação existe, mas é desigual. As políticas avançaram no plano
normativo; as práticas pedagógicas dispõem de um repertório relevante de
estratégias; a formação docente, no entanto, permanece como elo frágil,
sobretudo quando é concebida como responsabilidade individual e não como
política institucional.

Os resultados indicam que a inclusão não se realiza pela matrícula. A
presença física do estudante na sala regular é condição necessária, mas
não suficiente. A literatura analisada mostra que a participação e a
aprendizagem dependem de adaptação curricular, comunicação acessível,
mediação social, avaliação contínua, trabalho colaborativo e relação
família-escola. Mostra também que essas condições não são espontâneas:
elas exigem tempo, recursos, formação e decisão política. Quando faltam,
a inclusão tende a se converter em experiência de exclusão dentro da
escola.

O estudo contribui para o campo ao organizar a literatura em três eixos
e ao propor quatro tensões transversais --- norma e prática,
identificação e estigma, adaptação e segregação, formação inicial e
continuada. Essas tensões podem funcionar como instrumento de análise
para gestores e formadores, ajudando-os a identificar em que ponto uma
prática inclusiva se desvia de seus objetivos. A contribuição é
analítica, não prescritiva: o trabalho não oferece um modelo a ser
aplicado, mas um conjunto de distinções que podem qualificar a
discussão.

As limitações devem ser explicitadas. Trata-se de revisão narrativa, sem
busca sistemática, sem avaliação padronizada de risco de viés e sem
análise meta-analítica. O corpus é restrito à literatura nacional e a um
recorte temporal específico. Não foram realizadas observações de campo
nem entrevistas. Por isso, os resultados não descrevem a realidade de
todas as escolas brasileiras; indicam tendências e tensões que a
literatura permite sustentar. A verificação de metadados não substitui a
leitura crítica das obras nem a avaliação de sua qualidade metodológica.

Por fim, cabe reafirmar que a inclusão escolar é um direito, não uma
concessão. A escola que acolhe a diversidade não apenas cumpre a lei;
ela enriquece o processo educativo para todos os envolvidos. A presença
de estudantes com TEA nos anos iniciais não é um problema a ser
administrado, mas uma oportunidade de repensar o currículo, a avaliação,
a formação e a organização do trabalho pedagógico. A inclusão, nesse
sentido, é uma tarefa coletiva e contínua, que exige vontade política,
condições institucionais e compromisso ético com a aprendizagem de
todos.

\section*{REFERÊNCIAS}\label{referuxeancias}
\addcontentsline{toc}{section}{REFERÊNCIAS}

BEZERRA, Marli de Figueiredo; PANTONI, Rodrigo Palucci. Formação docente
para inclusão de alunos com Transtorno do Espectro Autista no Ensino
Médio Integrado. \textbf{Educitec -- Revista de Estudos e Pesquisas
sobre Ensino Tecnológico}, Manaus, v. 8, p.~e182622, 2022. DOI:
\url{10.31417/educitec.v8.1826}.

BRASIL. \textbf{Constituição da República Federativa do Brasil de 1988}.
Brasília, DF: Senado Federal, 1988. Disponível em:
\url{http://www.planalto.gov.br/ccivil_03/constituicao/constituicao.htm}.
Acesso em: 5 nov. 2026.

BRASIL. \textbf{Lei nº 9.394, de 20 de dezembro de 1996}. Estabelece as
diretrizes e bases da educação nacional. Brasília, DF: Presidência da
República, 1996. Disponível em:
\url{http://www.planalto.gov.br/ccivil_03/leis/l9394.htm}. Acesso em: 5 nov.
2026.

BRASIL. \textbf{Lei nº 12.764, de 27 de dezembro de 2012}. Institui a
Política Nacional de Proteção dos Direitos da Pessoa com Transtorno do
Espectro Autista. Brasília, DF: Presidência da República, 2012.
Disponível em:
\url{http://www.planalto.gov.br/ccivil_03/_ato2011-2014/2012/lei/l12764.htm}.
Acesso em: 5 nov. 2026.

BRASIL. \textbf{Lei nº 13.146, de 6 de julho de 2015}. Institui a Lei
Brasileira de Inclusão da Pessoa com Deficiência. Brasília, DF:
Presidência da República, 2015. Disponível em:
\url{http://www.planalto.gov.br/ccivil_03/_ato2015-2018/2015/lei/l13146.htm}.
Acesso em: 5 nov. 2026.

BRASIL. Ministério da Educação. \textbf{Política Nacional de Educação
Especial na Perspectiva da Educação Inclusiva}. Brasília, DF: MEC, 2008.
Disponível em: \url{http://portal.mec.gov.br/seesp/arquivos/pdf/politica.pdf}.
Acesso em: 5 nov. 2026.

CABRAL, Cristiane Soares; MARIN, Angela Helena. Inclusão escolar de
crianças com transtorno do espectro autista: uma revisão sistemática da
literatura. \textbf{Educação em Revista}, Belo Horizonte, v. 33,
p.~e142079, 2017. DOI: \url{10.1590/0102-4698142079}.

COUTO, Cirleine Costa; FURTADO, Maria Cândida de Carvalho; ZILLY,
Adriana; SILVA, Marta Angélica Iossi. Experiências de professores com o
autismo: impacto no diagnóstico precoce e na inclusão escolar.
\textbf{Revista Eletrônica de Enfermagem}, Goiânia, v. 21, p.~e55954,
2019. DOI: \url{10.5216/ree.v21.55954}.

FAM, Antônio Eustáquio de Oliveira; REIS, Sérgio Pereira dos; BARBOSA,
Rosemary Pereira Costa e. O brincar no espaço escolar como estratégia de
inclusão de alunos com Transtorno do Espectro Autista (TEA).
\textbf{Research, Society and Development}, {[}s. l.{]}, v. 10, n.~6,
p.~e49010615912, 2021. DOI: \url{10.33448/rsd-v10i6.15912}.

MAIA, Juliana; BATAGLION, Giandra Anceski; MAZO, Janice Zarpellon.
Alunos com transtorno do espectro autista na escola regular: relatos de
professores de educação física. \textbf{Revista da Associação Brasileira
de Atividade Motora Adaptada}, Marília, v. 21, n.~1, 2020. DOI:
\url{10.36311/2674-8681.2020.v21n1.02.p15}.

NUNES, Débora Regina de Paula; AZEVEDO, Mariana Queiroz Orrico de;
SCHMIDT, Carlo. Inclusão educacional de pessoas com autismo no Brasil:
uma revisão da literatura. \textbf{Revista Educação Especial}, Santa
Maria, v. 26, n.~47, p.~557-572, set./dez. 2013. DOI:
\url{10.5902/1984686X10178}.

SCHMIDT, Carlo et al.~Inclusão escolar e autismo: uma análise da
percepção docente e práticas pedagógicas. \textbf{Psicologia: Teoria e
Prática}, São Paulo, v. 18, n.~1, p.~222-235, 2016. DOI:
\url{10.15348/1980-6906/psicologia.v18n1p222-235}.

TOGASHI, Cláudia Miharu; WALTER, Cátia Crivelenti de Figueiredo. As
contribuições do uso da comunicação alternativa no processo de inclusão
escolar de um aluno com transtorno do espectro do autismo.
\textbf{Revista Brasileira de Educação Especial}, Marília, v. 22, n.~3,
p.~351-366, 2016. DOI: \url{10.1590/S1413-65382216000300004}.
